\subsection{Conclusion}

Having worked through both Aristotle's analyses and Heidegger's reading of them, the chapter's central claim can now be stated in its most concentrated form. Each of the higher-order \textit{pathē} that Aristotle catalogs in the \textit{Rhetoric} --- anger, calm, friendship, enmity, fear, confidence, shame, kindness, pity, indignation, envy, emulation --- is a determinate concretion of a single five-fold ontological structure. The five-fold is neither the three-fold of the \textit{Rhetoric} (change, \textit{krisis}, hedonic tonality) nor the four-fold of the \textit{Metaphysics} (alterability, actual alteration, harmful alteration, magnitude) nor the four-component composite of the enmattered-account analysis (propositional evaluation, conative orientation, hedonic tonality, physiological alteration); it integrates the three-fold and four-fold while making explicit the temporal and disclosive moments that the chapter has developed across \S\S1.4--1.8. Setting the structure out cleanly, then turning back to its grounding in the disclosive structure of \textit{Befindlichkeit}, completes the chapter's argument and orients the reader toward what follows.

\begin{adjustwidth}{0.5in}{0in}
\begin{enumerate}
\item \textbf{The disposition.} Each \textit{pathos} is a determinate manner of \textit{finding-oneself-disposed}. It is not an inner state that the human being inspects; it is the way the human being is taken in this moment, in this situation. Anger is finding-oneself-as-having-been-slighted-by-a-particular-person-who-could-have-acted-otherwise; shame is finding-oneself-as-having-fallen-short-before-an-audience-that-matters; pity is finding-oneself-as-confronted-with-undeserved-suffering-of-one-with-whom-I-could-imagine-myself. The dispositional moment is what \S1.2's articulational-concretion thesis recovered as the \textit{Befindlichkeit}-side of the \textit{pathos}-structure, and what \S1.4 specified as the involuntary uptake through which \textit{doxa} concretizes the basic affective valence into determinate emotion.

\item \textbf{The toward-which and the toward-whom.} Each \textit{pathos} has a determinate intentional structure: it is \textit{about} this and \textit{with respect to} those. The toward-which articulates the world-side of the disposition (what is at issue); the toward-whom articulates its social-ontological side (the persons to whom and before whom it is articulated). Anger is \textit{with respect to} the slight, \textit{toward} the slighter; pity is \textit{with respect to} undeserved harm, \textit{toward} the one suffering; shame is \textit{with respect to} the fall, \textit{before} the audience that matters. The two together articulate the aspectual disclosure that the \textit{pathos} accomplishes: the way the situation is given as this-and-that, for-me. \S1.4's per-emotion analyses of the \textit{Rhetoric}'s definitions (\textit{Rhetoric} II.2, 1378a30--b9 for anger; II.5, 1382a22--32 for fear; II.8, 1385b13--19 for pity; 1383b15--17 for shame) supplied the specific evaluative contents that fill out this intentional structure for each emotion in the catalog.

\item \textbf{The bodily form.} Each \textit{pathos} has a determinate bodily articulation that is constitutive, not symptomatic. The hot rush of anger, the cold of fear, the hot flush of shame, the welling of the eyes in pity --- these are part of the way the \textit{pathos} \textit{is} in the world. To strip a \textit{pathos} of its bodily form is to lose the phenomenon. The bodily form is \textit{logos}-articulated: it has the \textit{eidos} of the slight, the threat, the fall, but it is enmattered, meaning-bodily-bodied-forth. This is the moment that \S1.3's analysis of \textit{De Anima} I.1, 403a25--b19 secured: anger is the desire for retaliation realized in the boiling of blood about the heart, where neither the formal nor the material account is sufficient on its own. \S1.5 specified the same moment from the action-theoretic side: the bodily form \textit{is} the somatic preparation that translates desire into bodily readiness in the \textit{De Motu Animalium} causal chain (701b33--702a3).

\item \textbf{The futural-temporal moment.} Each \textit{pathos} is structured by a determinate temporality. Fear is futural --- the coming-toward of an anticipated evil. Anger is present-recoil-from-having-been-slighted, with a prospective vengeance. Shame is present-recoil-from-having-publicly-fallen-short, with an indeterminate future of being-thus-marked. Pity has a present-confrontation with another's present-suffering, a backward look at the events that produced it, and a futural-anticipation that it might or could befall the witness. Each \textit{pathos}, accordingly, articulates being-there in its three temporal ecstases --- having-been, present, projected-futural --- with a characteristic profile that distinguishes one \textit{pathos} from another. \S1.7's analysis of the synchronic chain and diachronic saturation specified the temporal structure beneath this characteristic profile: each present \textit{pathos} is always already shaped by the having-been of prior actualizations sedimented as dispositional ground, and Heidegger's \textit{Geworfenheit} names at the existential level what Aristotle's \textit{hexis} names at the perceptual-and-ethical level.

\item \textbf{The disclosive function: bringing to \textit{krisis}, bringing to \textit{logos}.} Each \textit{pathos} shifts the hearer toward taking-a-position. Each opens or closes specific possibilities of \textit{krisis}, of view-construction or decision. Each cultivates a disposition that, at its limit, brings to speech: anger brings the slighted to confront, fear brings the threatened to confer, shame brings the fallen-short to defend or reform, pity brings the witness to alleviate. The cultivation of the \textit{pathos} is the cultivation of the bringing-to-discourse, and rhetoric, on Heidegger's reading, is the systematic art of just this cultivation. \S1.8's analysis of the chain at $M_3 \rightarrow M_4$ specified the mechanism: emotion converts \textit{orexis} into a determinate form of desire under a specific evaluative description, and this determinate desire is precisely what carries the agent into the world-engaging act through which the \textit{pathos} discloses itself in speech and action.
\end{enumerate}
\end{adjustwidth}

These five elements together articulate the generic structure of the \textit{pathē}. Each Aristotelian emotion in the \textit{Rhetoric} is a specific filling-out of the five --- an articulational concretion in which the disposition, the world-side and social-side intentionality, the bodily form, the temporality, and the disclosive function are just so, and not otherwise. The catalog of the \textit{Rhetoric} is therefore not a list of psychological types but a phenomenology of the determinate ways in which a finite, embodied, speaking being finds itself disposed before the matters that concern it.

The five-fold structure has its philosophical center of gravity, accordingly, in the claim that the \textit{pathē} are the ground of \textit{logos}. The point is Heidegger's most concentrated formulation of what the GA 18 lectures recover from Aristotle, and the conclusion's argument turns on it:

\begin{adjustwidth}{0.5in}{0in}
The \gk{πάθη} (\textit{pathē}) are not merely an annex of psychical processes, but are rather \textit{the ground out of which speaking arises, and which what is expressed grows back into}; the \gk{πάθη}, for their part, are \textit{the basic possibilities in which being-there itself is primarily oriented toward itself}, finds itself. The primary being-oriented, the illumination of its being-in-the-world is not a \textit{knowing}, but rather a \textit{finding-oneself} that can be determined differently, according to the mode of being-there of a being. (BCAP, p. 176)
\end{adjustwidth}

Speaking, on this reading, is not a master operation that draws on the \textit{pathē} as resources; speaking arises out of the \textit{pathē} and returns to them, the feedback structure of which \S1.7 has specified at the \textit{pathos} level itself. The rhetorical art, accordingly, is the cultivation of the dispositional ground that brings being-there to discourse, and Aristotle's \textit{Rhetoric} is therefore not a manual of persuasion-techniques but a phenomenology of the disclosive structures through which being-there finds itself in the \textit{polis}. What Heidegger's recovery discloses is that the \textit{pathē} are not psychological phenomena to be added to a prior account of cognition; they are the prior. They are the modes in which being-there has its world disclosed to it, and on the basis of which both cognition and speaking become possible. To analyze anger, fear, shame, pity, indignation, kindness --- the family of higher-order emotions Aristotle catalogs --- is, on this reading, to undertake the phenomenology of disclosedness itself: the inquiry into how a finite, embodied, speaking being finds itself in the world before any explicit knowing or saying, and on the basis of which knowing and saying become possible at all. Heidegger's retrospective citation of the \textit{Rhetoric} as ``the first systematic hermeneutic of the everydayness of Being-with-one-another'' (SZ \S29, H.138; Eng. 178) is, in this light, not a passing historical note but a precise diagnostic: the \textit{Rhetoric}'s catalog of \textit{pathē} is the first sustained phenomenology of mood, and the GA 18 lectures are the preparatory work in which that phenomenology is recovered for the \textit{Befindlichkeit}-analytic of \textit{Being and Time}.

The \textit{pathē}, moreover, are co-original with \textit{hexeis}. Both are fundamental concepts of being, articulating the being-structure of beings whose way of being is ``\textbf{******}'' (BCAP, p. \textbf{******}).\footnote{Verbatim text and page reference to be supplied manually for the BCAP / GA 18 passage on \textit{pathē} and \textit{hexeis} as fundamental concepts of speaking-with-one-another.} Each \textit{pathos} has its possible \textit{hexis}-correlate: courage as the \textit{hexis} of fear, magnanimity as the \textit{hexis} of anger, friendliness as the \textit{hexis} of kindness; \textit{aretē}, in the most general sense, is the cultivation of the \textit{hexis} that fits the \textit{kairos}. \S1.7's analysis of the diachronic register of the feedback loop specified the temporal mechanism: \textit{hexeis} are what consolidate over time through the repeated actualization of \textit{pathos}-chains, so that the present \textit{pathos} is always shaped by the dispositional ground that prior actualizations have sedimented. The five-fold structure of any given \textit{pathos} is therefore not a snapshot but a moment within a temporally extended cultivation; the \textit{pathē} are the hinge between the singular event of being-affected and the durable ethical character that consolidates from many such events.

The disclosive role that the \textit{pathē} discharge in the perception-to-action chain depends on a feature the chapter has named at several points but that warrants explicit final statement: emotion responds to appearance, not to corrected judgment. If \textit{phantasia} supplies the appearance, \textit{doxa} is the moment of apophantic assertion, the unfolding from taking-as-something to taking-something-as-true; the \textit{pathos} that follows is a being-moved keyed to the appearance, not to whatever subsequent rational correction the agent may apply. This is precisely what makes the chain run with the speed and inevitability that Aristotle ascribes to it (\textit{De Anima} III.3, 427b21--24). Three illustrative phenomena confirm the analysis. The sound mistaken for a rattlesnake produces real fear, even though the corrected judgment will dissolve the mistake; false beliefs generate real emotions because emotion is keyed to the \textit{phantasmatic} appearance, not to its eventual correction. Art produces emotion, because the \textit{phantasmatic} representation suffices to occasion the affective movement, even when the audience knows the represented event is unreal. And emotion can contradict subsequent judgment, as when knowing that the rope is not a snake fails to neutralize the startle, because the rational correction does not unwind the prior \textit{phantasmatic}-doxastic uptake on which the affective movement was built.

Because \textit{phantasia} generates appearances beyond the immediately present, the temporal depth of emotional life follows directly. Aristotle states in \textit{De Memoria} that we remember when ``the movement caused by the object is preserved internally'' (\textit{De Memoria} 450a20--25). To remember grief is to feel grief again; to anticipate pleasure is to feel pleasure now in relation to a not-yet. When Aristotle writes in the \textit{Rhetoric} that ``both the man who remembers and the man who hopes will be attended by an imagination of what he remembers or hopes'' (\textit{Rhetoric} I.11, 1370a), he articulates the temporal depth of emotional life directly: emotion is not anchored to the present but reaches into past and future by means of the \textit{phantasmatic} appearance. Fear, paradigmatically, is ``a pain or disturbance due to imagining some destructive or painful evil in the future'' (\textit{Rhetoric} II.5, 1382a22--23); the feared event does not exist, the \textit{phantasma} does, and the soul moves as though the feared thing were present. \textit{Phantasia} is therefore not a decorative adjunct to emotion; it is the condition of possibility for the soul's being-moved by absences, possibilities, appearances, and counterfactuals --- which is to say, the condition of possibility for the entire range of higher-order \textit{pathē} that the chapter's analysis has placed at the center of evaluatively complex action.

The practical, world-shaping power of emotion that Aristotle's rhetorical theory makes explicit follows from this analysis. To persuade is not primarily to offer proofs but to reshape the listener's \textit{phantasmatic} field. The rhetor must supply images --- concrete, narrative, affectively charged --- that allow the audience to see what the rhetor wants them to see, since this imagistic seeing grounds the subsequent formation of \textit{doxa}, the ensuing movement of emotion, and the eventual act of judgment. If rhetorical persuasion succeeds, accordingly, it is because it sets off the entire cascade of movement: new \textit{phantasma}, new \textit{doxa}, new emotion, new orientation, new action. Heidegger's gloss on this practical-philosophical fact in the GA 18 lectures and the \textit{Heidegger and Rhetoric} editorial reception captures the stakes precisely: rhetoric is the discipline in which ``the self-elaboration of Dasein is expressly executed'' (\textit{Heidegger and Rhetoric}, ed. Gross, p. 1).

For Aristotle, in sum, emotion is the kinetic, affective, evaluative dimension of how the world appears to a finite, embodied, speaking being. It is grounded in \textit{phantasia}, shaped through involuntary \textit{doxa}, felt as pleasure or pain, prepared somatically, and expressed as motion toward or away from what the soul takes to be significant. Emotion is not an add-on to rationality, nor a disruption of cognition, nor a mere physiological disturbance, but the way the soul's self-movement discloses significance at all. The five-fold structure --- disposition, intentionality, bodily form, futural-temporal moment, disclosive function --- specifies the architecture of this disclosure for each of the determinate \textit{pathē} the \textit{Rhetoric} catalogs. Emotion, in its full Aristotelian sense, is the affective architecture of being-in-the-world.

What the chapter has established opens specific questions that the rest of the work will pursue. The cultivation of the \textit{hexeis} that correspond to each \textit{pathos} --- the practical-ethical task of shaping one's dispositional ground over time --- is the subject of the next chapter, where the temporal sedimentation specified in \S1.7 receives its sustained practical-philosophical treatment. The \textit{kairos}-relativity of the appropriate \textit{pathos}, that is, the question of which emotion fits which situation and how the agent learns to recognize the fit, follows naturally from the \textit{hexis}-analysis: cultivation is always cultivation of the disposition that responds rightly to the right occasion at the right time. And the rhetorical case-studies through which the chapter's general apparatus is brought to bear on specific public-philosophical problems --- the case-studies in which the five-fold structure does its concrete diagnostic work --- will occupy the chapters thereafter. With the structural ground secured, accordingly, the analysis can now turn to the practical and rhetorical work that the structure makes possible, in which the architecture this chapter has specified does its world-engaging labor.
